%!TEX root = ../../main.tex

\chapter{Klassifikation der Prüfzifferverfahren}
\label{ch:klassifikation}

Die Welt der Prüfzifferverfahren ist vielfältig und facettenreich. Um einen strukturierten Überblick zu ermöglichen, werden diese Verfahren nach ihren mathematischen Grundprinzipien klassifiziert. Diese Klassifikation erleichtert nicht nur das Verständnis der Verfahren, sondern ermöglicht auch eine gezielte Auswahl je nach Anwendungsfall und Anforderungsprofil.

Im Folgenden werden vier Hauptkategorien von Prüfzifferverfahren vorgestellt, die sich in ihren mathematischen Grundlagen, ihrer Komplexität und ihren Fehlererkennungseigenschaften unterscheiden. Einfache Modulo-Verfahren nutzen die elementare Restklassenarithmetik für die Fehlererkennung. Sie sind leicht zu implementieren und zu verstehen, weshalb sie besonders in einfachen Anwendungsfällen zum Einsatz kommen, wie etwa bei älteren Bankleitzahlen Gewichtsverfahren erweitern das Konzept durch positionsabhängige Gewichtungsfaktoren und erhöhen dadurch die Fehlererkennungsleistung deutlich. Sie finden Anwendung in wichtigen Identifikationssystemen wie der ISBN für Bücher oder den EAN-Codes im Einzelhandel. Verfahren mit Permutation führen zusätzlich eine systematische Umordnung der Ziffern durch, wodurch komplexere Fehlertypen erkannt werden können. Der Luhn-Algorithmus für Kreditkartennummern ist hierfür ein Paradebeispiel. Polynomielle Verfahren schließlich repräsentieren die mathematisch anspruchsvollste Kategorie, basierend auf der Arithmetik von Polynomen über endlichen Körpern. Der CRC (Cyclic Redundancy Check) ist der wichtigste Vertreter dieser Kategorie und wird häufig in der digitalen Datenübertragung eingesetzt.

Für jede dieser Kategorien werden in den folgenden Abschnitten die mathematischen Grundlagen erläutert und anhand von konkreten Beispielen veranschaulicht. Dabei wird sowohl auf die theoretischen Aspekte als auch auf die praktische Anwendung eingegangen, um ein umfassendes Verständnis zu vermitteln \cite{Hoffmann2023-sv}.
